\subsection{关于双模的两个引理}

设 A 与 B 是环。对从 B 到 A 的任一同态 f，我们用 \(A^{f}\) 表示以右 A-模 \(A_{d}\) 为基础、其左 B-模结构的外部法则由 \((b,a)\mapsto f(b)a\) 给出的 (B,A)-双模。

引理 1. —— 设 f 与 g 是从 B 到 A 的同态。下述条件等价：

\begin{enumerate}
\def\labelenumi{(\roman{enumi})}
\item
  (B,A)-双模 \(A^{f}\) 与 \(A^{g}\) 同构。
\item
  存在 A 的一个内自同构 \(\theta\)（I, §8, No.~4, 第 102 页, 例 2），使得 \(g = \theta \circ f\)。
\end{enumerate}

右 A-模 \(A_{d}\) 的自同构是映射 \(x \mapsto ax\)，其中 a 是 A 的一个可逆元。这样的一个自同构是从 \(A^{f}\) 到 \(A^{g}\) 的 B-线性映射当且仅当我们有

\[
g (b) a x = a f (b) x
\]

对 A 中的每个 x 与 B 中的每个 b。这个关系式等价于对 B 中的每个 b 有 \(g(b) = af(b)a^{-1}\)，即等价于 \(g = \theta \circ f\)，其中 \(\theta\) 是 A 的内自同构 \(x \mapsto axa^{-1}\)。

引理 2. —— 假设 B 是一个半单环，且是其中心 Z 上的一个有限生成模。设 M 与 N 是 (B, A)-双模。假设它们有有限长度（特别地，当它们是有限长度的右 A-模时就是这种情形）。若 M 与 N 作为 (Z, A)-双模同构，则它们作为 (B, A)-双模同构。

\begin{enumerate}
\def\labelenumi{\Alph{enumi})}
\item
  首先考虑 B 是交换域 L 上一个有限维 d 向量空间 S 的自同态环的情形。这时 Z = L；我们把 S 看作一个 (B, Z)-双模。环 B 是单的，S 是一个单 B-模，而 Z 是 S 的交换子；每个 B-模都是 S 型的同型模（VIII, 第 122 页, 命题 2 a)）。设 \((V, \alpha)\)（相应地 \((W, \beta)\)）是 B-模 M（相应地 N）的一个描述。集合 V（相应地 W）带有一个 (Z, A)-双模结构，使得 \(\alpha\)（相应地 \(\beta\)）是 (B, A)-双模的同构（VIII, 第 64 页, 注 2）。作为 (Z, A)-双模，M 同构于 \(V^{d}\) 而 N 同构于 \(W^{d}\)，且存在一个从 V 的 (Z, A)-子双模的集合（按包含排序）到 M 的 (B, A)-子双模的集合的同构（同前所引）。于是 V 是一个有限长度的 (Z, A)-双模，W 亦然。由于 (Z, A)-双模 \(V^{d}\) 与 \(W^{d}\) 同构，由应用于环 \(Z \otimes_{Z} A^{\circ}\) 的 VIII, 第 37 页的定理 2, d)，(Z, A)-双模 V 与 W 同构。最后，(B, A)-双模 M 与 N 同构。
\item
  现在考虑 B 是一个作为 Z-模有限生成的单环的情形。这时 Z 是一个域，而 B 是域 Z 上的一个有限次数的中心单代数。由 VIII, 第 252 页的定理 1，存在 Z 的一个在 Z 上次数有限的扩张 \(Z'\)，使得 \(Z'\) -代数 \(B' = B_{(Z')}\) 同构于一个有限维 \(Z'\) -向量空间的自同态代数。令 \(M' = M_{(Z')}\) 与 \(N' = N_{(Z')}\)。则 \(M'\) 与 \(N'\) 是有限长度的 \((B', A)\) -双模；看作 \((Z', A)\) -双模时，\(M'\) 与 \(N'\) 同构。由 A) 中处理的情形，\(M'\) 与 \(N'\) 作为 \((B', A)\) -双模同构，从而更不用说作为 \((B, A)\) -双模同构。令 \(r = [Z' : Z]\)。\((B, A)\) -双模 \(M' = Z' \otimes_{Z} M\) 同构于 \(M^r\)，同样地，\((B, A)\) -双模 \(N'\) 同构于 \(N^r\)。由于 M 与 N 是有限长度的 \((B, A)\) -双模，由 VIII, 第 37 页的定理 2, d) 可知 \((B, A)\) -双模 M 与 N 同构。
\item
  最后考虑一般情形，即 B 是一个作为 Z-模有限生成的半单环。设 S 是单 B-模的类所成的集合；它是有限的（VIII, 第 136 页, 命题 1）。对任何 \(\lambda \in S\)，用 \(M_{\lambda}\)（相应地 \(N_{\lambda}\)）表示 B-模 M（相应地 N）的 \(\lambda\) 型同型分量；它是 M（相应地 N）的一个 (B, A)-子双模（VIII, 第 67 页的注）。对 \(\lambda \in S\)，把 B-模 \(\lambda\) 的零化子记作 \(b_{\lambda}\)，并令 \(B_{\lambda} = B / b_{\lambda}\)；设 \(Z_{\lambda}\) 是 \(B_{\lambda}\) 的中心。对 \(\lambda \in S\)，\((B_{\lambda}, A)\) -双模 \(M_{\lambda}\) 与 \(N_{\lambda}\) 有有限长度。于是我们可以把 B 等同于单环 \(B_{\lambda}\) 的乘积，把 Z 等同于 \(Z_{\lambda}\) 的乘积（VIII, 第 141 页, 命题 8）。此外，我们可以把 M 等同于 \(\prod_{\lambda \in S} M_{\lambda}\)，把 N 等同于 \(\prod_{\lambda \in S} N_{\lambda}\)。由假设，M 与 N 作为 (Z, A)-双模同构；由此得出，对 \(\lambda \in S\)，\(M_{\lambda}\) 与 \(N_{\lambda}\) 作为 ( \(Z_{\lambda}, A\) )-双模同构。由 B) 中处理的情形，\((B_{\lambda}, A)\) -双模 \(M_{\lambda}\) 与 \(N_{\lambda}\) 同构，于是 (B, A)-双模 M 与 N 同构。
\end{enumerate}

注. —— 由引理 2 的证明可知 M 与 N 是有限长度的 (Z, A)-双模。因此，若 B 与 A 是两个作为其各自的中心 Z(B) 与 Z(A) 上的有限生成模的半单环，则两个作为 (Z(B), Z(A))-双模同构的有限长度 (B, A)-双模是同构的。
% semantic-fragments: legacy-input-seam
\relax{}
